\section{Reverse-complement molecules and bidirected flows}
\label{app:molecular-flow}

The preceding oriented word results do not automatically apply to
reverse-complement molecule classes or to disconnected bidirected flows.
This appendix records the source-faithful same-length counterexamples,
their essential distinctions, and a compact comparison of model domains.
All secondary flow configurations, exact incidence matrices, detailed
likelihood factors and the earlier exhaustive model table are preserved
in \path{paper/technical-notes/full-molecular-flow.tex} and
\path{docs/verification.md}.

\begin{theorem}[MB09 variable-length molecule counterexample]
\label{thm:cex-molecule}
In the DNA-molecule model with $A\leftrightarrow T$ and $L=3$,
the truth $S=\texttt{AAATT}$ with observed starts $(0,1,4)$
and the longer competitor $D=\texttt{AAAATT}$ both induce
MB09-admissible bidirected circuits. With the fixed external
reference $N=5$, the original Section~6.1 surrogate gives
$\Lik_{6.1}(D)/\Lik_{6.1}(S)=9/8>1$.
The candidate-intrinsic exact multinomial likelihood also
favors $D$, with ratio $125/108>1$.
\end{theorem}
\begin{proof}
The observed molecule counts are $(1,1,1)$ on
$(\texttt{AAA},\texttt{AAT},\texttt{TAA})$; the candidate
counts are $(1,2,2)$ for $S$ and $(2,2,2)$ for $D$.
Only the first fixed-$N$ binomial factor changes:
\[
\frac{(2/5)(3/5)^2}{(1/5)(4/5)^2}=\frac98.
\]
For exact likelihood, the observed molecule-class probabilities
are $(1/5,2/5,2/5)$ under $S$ and
$(2/6,2/6,2/6)$ under $D$, giving
\[
\frac{(2/6)^3}{(1/5)(2/5)^2}=\frac{125}{108}>1.
\]
The multinomial coefficient cancels. The physical circuit feasibility is verified by the exact
source-model certificate indexed in Appendix~\ref{app:reproduction}.
\end{proof}

\subsection{The historical matched-string predicate}
\label{sec:finite:historical}

The supplement to \textcite{shomorony2016} formulates coverage
and repeat bridging in terms of \emph{observed read strings matched
anywhere on the truth}, rather than solely their sampled latent
placements. The difference matters whenever a string occurs more
than once. The definitions are given in
Definition~\ref{def:historical-match}.

\begin{proposition}[Historical matching and latent placement are incomparable]
\label{prop:historical-source}
Neither the historical read-string predicate
$I_s^{\mathrm{match}}$ nor the sampled-placement predicate
$\Is$ implies the other in general.
\end{proposition}
\begin{proof}
On the circle $\texttt{AAAA}$ with $L=3$, one observed
$\texttt{AAA}$ matches at every start and meets the
historical coverage/bridging clauses, but its sole sampled
placement at start $0$ leaves one base uncovered.
Conversely, on $\texttt{ACGT}$ with $L=2$, placements at
starts $0,2$ cover every base, but the historical coverage
clause requires read strings matching also starts $1,3$,
which are absent.
\end{proof}

\begin{theorem}[Historical same-length molecule counterexample $W_1$]
\label{thm:cex-historical-W1}
With $A\leftrightarrow T$, $L=3$, truth
$S=\texttt{AAATAT}$ of length $6$ and observed starts
$(0,0,1,3,5)$, the candidate $D=\texttt{AAAAAT}$
has the same length. The truth satisfies historical
matched-string coverage and every literal triple
and interleaved-repeat bridging clause; both genomes
are MB09-feasible full-overlap bidirected circuits.
The exact molecule likelihood ratio is $3$, and the
fixed-$N$ binomial-surrogate ratio is $5$, in favour
of $D$.
\end{theorem}
\begin{proof}[Finite certificate]
In class order
$(\texttt{AAA},\texttt{AAT},\texttt{ATA/TAT},\texttt{TAA})$,
the multiplicities are
\[
  x=(2,1,1,1),\quad c_S=(1,1,3,1),\quad c_D=(3,1,1,1).
\]
The ratios follow from the changed first and third
factors; each observed class has positive candidate
count. The full bridging and physical-flow certificates are
checked in the Lean module indexed in
Appendix~\ref{app:reproduction}.
\end{proof}

\begin{theorem}[Historical same-length molecule counterexample $W_2$]
\label{thm:cex-historical-W2}
On the four-letter DNA alphabet with usual reverse
complement, let $L=3$, $S=\texttt{ATATACAC}$ and
$D=\texttt{ATACACAC}$, both of length $8$, with
observed starts $(1,3,4,5,6,7)$.
The historical matched-string predicate holds at full
strength, and both genomes are MB09-admissible circuits.
Moreover, both satisfy the \emph{additional}, non-source
per-occurrence inequality $x(w)\le d_S(w),d_D(w)$
for every observed molecule type. Nevertheless
\[
 \frac{\Lik_{\mathrm{ex}}(D;x)}{\Lik_{\mathrm{ex}}(S;x)}
 =\frac32>1,\qquad
 \frac{\Lik_{\mathrm{bin}}(D;x)}{\Lik_{\mathrm{bin}}(S;x)}
 =\frac95>1.
\]
\end{theorem}
\begin{proof}[Finite certificate]
In molecule-class order
$(\texttt{ATA/TAT},\texttt{TAC/GTA},
  \texttt{ACA/TGT},\texttt{CAC/GTG},\texttt{CAT/ATG})$,
the observed counts are $(1,1,2,1,1)$,
the truth's $(3,1,2,1,1)$ and the competitor's
$(1,1,3,2,1)$. Thus both dominate the observations
coordinatewise, while the exact product and fixed-$N$
product give the displayed ratios. The sampled placements
also satisfy the weaker placement predicate for this
instance, but the theorem uses the \emph{historical}
read-string certificate. Repeat bridging, bidirected incidence and feasibility
are checked by the historical-witness Lean modules
indexed in Appendix~\ref{app:reproduction}.
\end{proof}

The per-occurrence strengthening in $W_2$ is not the
original MB09 rule, which requires a lower bound on
\emph{vertex} occurrence instead. The same witness
therefore subsumes the earlier placement-only
per-occurrence example; we state it just once.
The witnesses also distinguish a spelled connected
circuit from an arbitrary balanced integral or
half-integral flow. Neither theorem is a claim about
the $2G$ artificial strand-concatenation model.

\subsection{Comparison by candidate universe}

\begin{longtable}{@{}p{0.34\textwidth}p{0.22\textwidth}p{0.32\textwidth}@{}}
\caption{Representative model outcomes; the complete case ledger
remains in the repository's technical note.}
\label{tab:finite}\\
\toprule
Model and condition & candidates & proved conclusion\\
\midrule
\endfirsthead
\toprule Model and condition & candidates & proved conclusion\\
\midrule\endhead
\bottomrule\endlastfoot
Oriented, realized placement bridging, same length and complete support
& circular spellings
& Rotation uniqueness by structural rigidity
(Thm.~\ref{thm:rigidity})\\
Oriented, historical matched-string bridging, same length and complete support
& circular spellings
& Rotation uniqueness (Cor.~\ref{cor:historical-oriented-same-length})\\
Oriented, exact iid reads, free length
& all circular genomes
& Strict failures (Thm.~\ref{thm:cex-variable-length});
sometimes no maximum (Thm.~\ref{thm:unattained-ML})\\
Oriented, primitive-P2 truth with $B<2G$
& all genomes of length at most $B$
& Exponential ML recovery
(Cor.~\ref{cor:length-window-ML})\\
MB09 molecules, historical matched-string bridging, same length
& spelled bidirected circuits
& Strict exact and surrogate failures $W_1$, $W_2$
(Thms.~\ref{thm:cex-historical-W1}--\ref{thm:cex-historical-W2})\\
MB09 variable-length molecules
& spelled bidirected circuits
& Strict exact and surrogate failures (Thm.~\ref{thm:cex-molecule})\\
MB09 general $o_{\min}$ flow optimization
& integral / relaxed bidirected flows
& General optimizer-to-spelling correspondence remains open\\
Two-strand $2G$ concatenation
& model of \textcite{bresler2013}
& Separate determinate model; general implication open\\
Disjoint two-circle duplex
& model-dependent
& Source interpretation differs from artificial concatenation
(Remark~\ref{rem:duplex})\\
\end{longtable}

The two historical same-length molecule failures settle those
\emph{specified} likelihood claims. Separate open questions
concern which exact objective and candidate universe the 2016
sentence intended, the general-$o_{\min}$ bidirected-flow
optimizer-to-spelling correspondence, and the different
two-strand concatenation model. They should not be conflated
into a single unresolved row. A complete source comparison
is retained in
\path{paper/technical-notes/full-molecular-flow.tex}.
